\chapter*{From the circular to Java}
\addcontentsline{toc}{chapter}{From the circular to Java}
\markboth{From the circular to Java}{}
\label{map:overview}
\begingroup
\tikzset{
 mapbox/.style={draw=teal,fill=pale,rounded corners=2pt,align=center,
   text width=42mm,minimum height=16mm,inner sep=3mm,font=\small},
 mapwide/.style={mapbox,text width=108mm,minimum height=11mm,inner sep=2mm},
 mapresearch/.style={mapbox,draw=navy,dashed,fill=white},
 maparr/.style={-{Latex[length=2mm]},thick,draw=navy},
 mapnote/.style={font=\small,align=center,fill=white,inner sep=1mm}
}
\newcommand{\MapSubheading}[1]{\par\vspace{2mm}{\sffamily\bfseries\large #1\par}\vspace{1mm}}
\newcommand{\MapHeading}[2]{\phantomsection\label{#2}\MapSubheading{#1}}
\newenvironment{MapDiagram}{\par\noindent\begin{minipage}{\linewidth}\centering}{\par\end{minipage}\par}
\newcommand{\MapCaption}[1]{\par{\small\itshape #1\par}\vspace{2mm}}

A circular supplies words; a Java program needs precise decisions. Between
them lie questions about meaning, evidence and execution. The eight steps
below locate that work. Each numbered box links to a closer view on the next
pages, where the technology is introduced through the job it performs.

\begin{MapDiagram}
\begin{tikzpicture}[node distance=5mm]
\node[mapwide] (a) {\hyperref[map:source]{\textbf{1. Retain the circular}}\\Original document and usable text};
\node[mapwide,below=of a] (b) {\hyperref[map:context]{\textbf{2. Read it in context}}\\Related sources, dates and business scope};
\node[mapwide,below=of b] (c) {\hyperref[map:readings]{\textbf{3. Develop possible interpretations}}\\Duties, exceptions and required facts};
\node[mapwide,below=of c] (d) {\hyperref[map:challenge]{\textbf{4. Challenge the interpretations}}\\Omissions, alternative readings and revealing cases};
\node[mapwide,below=of d] (e) {\hyperref[map:rules]{\textbf{5. Express the proposed rules precisely}}\\Separate duties and alternative meanings};
\node[mapwide,below=of e] (f) {\hyperref[map:java]{\textbf{6. Generate Java and check its behaviour}}\\Execution tests and scoped mathematical checks};
\node[mapwide,below=of f] (g) {\hyperref[map:delivery]{\textbf{7. Deliver a reviewable Java program}}\\Program, sources, results and open questions};
\node[mapwide,dashed,fill=white,below=of g] (h) {\hyperref[map:operation]{\textbf{8. Approve and supervise bank use}}\\Institution decisions, operating controls and source review};
\foreach \x/\y in {a/b,b/c,c/d,d/e,e/f,f/g,g/h}
  \draw[maparr] (\x)--(\y);
\draw[maparr] (d.east)--++(12mm,0)|- node[pos=.25,mapnote,rotate=90]{Revise} (c.east);
\draw[maparr] (f.west)--++(-12mm,0)|- node[pos=.25,mapnote,rotate=90]{Repair} (e.west);
\end{tikzpicture}
\end{MapDiagram}
\MapCaption{A tested program is a reviewable result. The dashed final step
requires the bank's own authority and operating evidence.}

The implemented route can produce draft controls while interpretations remain
unresolved. It can also stop before code generation when a proposal cannot be
expressed or evidence is missing. In the running example, whether a rebate is
a fee discount can remain open even after its proposed rule runs correctly.
The later charts show how that question stays attached to the result.

\clearpage
\MapHeading{1. Retain the circular}{map:source}
The first task is to give every later reading the same inspectable source.
For PDFs, \texttt{pypdf} and Poppler's \texttt{pdftotext} are two different text
extractors. Web pages have two parsing paths. Comparing their text can expose
a missing line; checking the page for images can expose material that neither
extractor has read.

\begin{MapDiagram}
\begin{tikzpicture}[node distance=7mm and 5mm]
\node[mapbox] (a) {Official circular\\PDF or web document};
\node[mapbox,right=of a] (b) {Keep the original\\Locate pages and passages};
\node[mapbox,right=of b] (c) {Extract the text\\Two extraction paths};
\node[mapbox,below=of b] (d) {Compare the extractions\\Inspect visual warnings};
\node[mapbox,below=of a] (e) {Retain usable text\\with its source locations};
\node[mapbox,below=of c] (f) {Unread or inconsistent?\\Keep the gap visible};
\draw[maparr] (a)--(b);\draw[maparr] (b)--(c);
\draw[maparr] (c.south)--++(0,-3mm)-|(d.north);
\draw[maparr] (d)--(e);
\draw[maparr] (d)--(f);
\end{tikzpicture}
\end{MapDiagram}
\MapCaption{Agreement between extractors helps find discrepancies; it does not
read a scanned exception. Automatic reading of scanned text is not implemented.}

\MapHeading{2. Read it in context}{map:context}
A gift restriction may refer to the Code of Conduct and a frequently asked
question (FAQ). A retained authority catalog connects those references to
document editions and passages. A dependency search follows explicit official
links; missing documents remain named questions. The relevant business scope
identifies the people, products, activities and dates being assessed.

\begin{MapDiagram}
\begin{tikzpicture}[node distance=7mm and 5mm]
\node[mapbox] (a) {Circular passage\\and cited references};
\node[mapbox,right=of a] (b) {Authority catalog\\Edition and passage lookup};
\node[mapbox,right=of b] (c) {Date and scope check\\What can this source support?};
\node[mapbox,below=of b] (d) {Source set for the review\\with missing or uncertain items};
\draw[maparr] (a)--(b);\draw[maparr] (b)--(c);
\draw[maparr] (c.south)--++(0,-3mm)-|(d.north);
\end{tikzpicture}
\end{MapDiagram}
\MapCaption{Finding the FAQ changes the question from where the source is to
what its dated answer supports.}

The retained FAQ page has later answer updates. Its older page date therefore
does not establish everything a historical reader could have seen. A catalog
match records useful evidence while preserving that uncertainty.
Chapter~\ref{ch:circular} follows the actual Code and FAQ example;
the companion's Appendix~\ref{app:current-interpretation} gives the source tools.

\clearpage
\MapHeading{3. Develop possible interpretations}{map:readings}
A language model can propose what a passage requires and which facts would
decide a case. Separate requests produce initial readings without showing each
request the other answers. Another set of requests inventories the source's
requirements and qualifications. This second view helps expose an exception
that all the proposed readings might miss.

\begin{MapDiagram}
\begin{tikzpicture}[node distance=7mm and 5mm]
\node[mapbox] (a) {Source passages\\Review question};
\node[mapbox,right=of a] (b) {Separate model requests\\Readings and facts};
\node[mapbox,right=of b] (c) {Alternative readings\\with source reasons};
\node[mapbox,below=of b] (d) {Source inventory\\Duties and qualifications};
\node[mapbox,below=of c] (e) {Questions to investigate\\Differences and possible gaps};
\draw[maparr] (a)--(b);\draw[maparr] (b)--(c);
\draw[maparr] (a.south)--++(0,-3mm)-|(d.north);
\draw[maparr] (d)--(e);\draw[maparr] (c)--(e);
\end{tikzpicture}
\end{MapDiagram}
\MapCaption{Several readings reveal disagreement. A separate source inventory
can reveal a shared omission. Separate requests can still share mistakes.}

\MapHeading{4. Challenge the interpretations}{map:challenge}
Tree search organises successive revisions of a reading. Breadth-first search
visits the shallower alternatives first. Upper Confidence bounds applied to
Trees (UCT) balances previously useful branches against less-visited ones.
These are ways to choose the next investigation, not probabilities of legal
correctness. Both are implemented with finite work limits.

\begin{MapDiagram}
\begin{tikzpicture}[node distance=8mm and 5mm]
\node[mapbox] (a) {Select a reading\\Breadth-first or UCT};
\node[mapbox,right=of a] (b) {Challenge a reading\\Sources, examples, cases};
\node[mapbox,right=of b] (c) {Record the consequence\\Revision, rejection or open question};
\node[mapbox,below=of b] (d) {Keep alternatives and evidence\\Stop at the work limit};
\draw[maparr] (a)--(b);\draw[maparr] (b)--(c);
\draw[maparr] (c.north)--++(0,5mm)-|node[pos=.25,above,mapnote]{further work justified}(a.north);
\draw[maparr] (c.south)--++(0,-3mm)-|(d.north);
\end{tikzpicture}
\end{MapDiagram}
\MapCaption{Search changes a proposed reading only with an explicit reason.
Exhausting the work allowance leaves an unresolved result.}

Typed arguments record support and objections; official examples retain their
stated meaning. A rebate-and-voucher case can reveal where readings attach an
exception. After aligning facts and outputs, the original programs execute
that case. Draft compilation and step~6 therefore occur inside this loop too.
Agreement still cannot establish complete source coverage.
Chapters~\ref{ch:search} and~\ref{ch:ensemble} develop these checks.

\clearpage
\MapHeading{5. Express the proposed rules precisely}{map:rules}
A program needs the meaning of each fact and operation fixed. LegalMath's rule
intermediate representation, \textbf{RuleIR}, is its small executable rule
language. It records typed facts, conditions, exceptions and scope. Type checks
reject such mistakes as adding a date to money; dependency checks reject
missing or circular rule references.

\begin{MapDiagram}
\begin{tikzpicture}[node distance=5mm]
\node[mapwide] (a) {Proposed interpretation\\Identify each duty, fact and exception};
\node[mapwide,below=of a] (b) {RuleIR expressions\\Precise operations with known, unknown and conflicting facts};
\node[mapwide,below=of b] (c) {Validate types and dependencies\\Unsupported expressions remain explicit gaps};
\node[mapwide,below=of c] (d) {Compose separately scoped controls\\Keep rival selections in separate draft packages};
\foreach \x/\y in {a/b,b/c,c/d}\draw[maparr](\x)--(\y);
\end{tikzpicture}
\end{MapDiagram}
\MapCaption{The current Java route starts from validated RuleIR. Successful
construction does not decide which legal interpretation the bank should adopt.}

For example, a gift restriction and a separately specified disclosure duty can
both belong in a package. Two rival readings of the restriction remain
alternative packages. The proposed grouping and each unresolved question are
preserved. Catala's treatment of definitions and exceptions informs the book's
design discussion; the current emitter does not invoke the Catala compiler.

\MapSubheading{Where Stipula belongs}
Stipula describes contracts through parties, resources, actions and state
changes. It is relevant when consent can be granted and withdrawn. The
published translation studied in Chapter~\ref{ch:languages} constructs a Java
verification model with specifications in the Java Modeling Language (JML).
KeY is a verifier that reasons about the Java model and those specifications.

\begin{MapDiagram}
\begin{tikzpicture}[node distance=5mm]
\node[mapresearch] (a) {Stipula contract\\Parties and state changes};
\node[mapresearch,right=of a] (b) {Java model with JML\\Conditions and invariants};
\node[mapresearch,right=of b] (c) {KeY verification\\Named properties of that model};
\draw[maparr,dashed](a)--(b);\draw[maparr,dashed](b)--(c);
\end{tikzpicture}
\end{MapDiagram}
\MapCaption{Dashed boxes show a studied research route, not the current
LegalMath compiler. Its Java model is not automatically a bank application.}

The retained translator examples contain a compilation discrepancy; no local
KeY replay establishes their reported proofs. Adoption requires reproducible
inputs and a justified mapping to LegalMath's event rules.
Chapter~\ref{ch:languages} details those limits; Chapter~\ref{ch:verification}
explains the implemented rule language.

\clearpage
\MapHeading{6. Generate Java and check its behaviour}{map:java}
The Java emitter embeds the fixed RuleIR rules in a generated policy class
that calls a shared Java evaluator. The Java compiler, \texttt{javac}, compiles
that class and its supporting runtime. A separately written Python evaluator
supplies another execution of the same proposed rule. Both receive the same
facts and assessment times.

\begin{MapDiagram}
\begin{tikzpicture}[node distance=9mm and 5mm]
\node[mapbox] (a) {Validated RuleIR\\Fixed proposed rule};
\node[mapbox,right=of a] (b) {Java emitter\\Policy class and shared evaluator};
\node[mapbox,right=of b] (c) {Java compiler\\Compiled program};
\node[mapbox,below=of a] (d) {Python evaluation\\Same facts and times};
\node[mapbox,below=of c] (e) {Java evaluation\\Same facts and times};
\node[mapbox,below=of d] (f) {Compare results\\Decision and explanation\\Required evidence};
\node[mapbox,right=of f] (g) {Keep test results\\Repair disagreements};
\draw[maparr](a)--(b);\draw[maparr](b)--(c);
\draw[maparr](a)--(d);\draw[maparr](c)--(e);\draw[maparr](d)--(f);
\draw[maparr](e.south)--++(0,-4mm)-|(f.north);
\draw[maparr](f)--(g);
\end{tikzpicture}
\end{MapDiagram}
\MapCaption{Two implementations check whether the proposed rule has survived
translation. They can agree on a rule that omits a legal exception.}

The test cases must also be capable of exposing mistakes. Boundary cases place
values on either side of a threshold. Missing-fact cases check that uncertainty
survives. Event histories exercise consent, withdrawal and timing. Deliberately
altered rules, called \emph{mutations}, ask whether the tests notice a removed
exception or a reversed condition.

\begin{MapDiagram}
\begin{tikzpicture}[node distance=5mm]
\node[mapbox] (a) {Introduce a defect\\Remove an exception};
\node[mapbox,right=of a] (b) {Try revealing cases\\Discount and voucher};
\node[mapbox,right=of b] (c) {Inspect the result\\Was the defect detected?};
\draw[maparr](a)--(b);\draw[maparr](b)--(c);
\end{tikzpicture}
\end{MapDiagram}
\MapCaption{A useful test notices a meaningful error. Mutation testing checks
that ability; it does not manufacture independent legal answers.}

For a small finite rule, every declared input combination can be enumerated.
Larger or unsupported domains require more limited conclusions. Expected legal
outcomes need their own source justification, separate from Python/Java
agreement. Compilation failures and disagreements return to the affected rule
or implementation; successful checks accompany the draft rather than remove
its open interpretation questions. Chapter~\ref{ch:verification} develops
these checks and the separate Java host demonstration.

\clearpage
\MapHeading{6 continued. Where mathematical verification fits}{map:formal}
A mathematical check needs a precise proposition: whether two encoded rules
differ on a permitted input, for example, or whether a transfer preserves a
quantity. Such propositions concern the formal model and its assumptions.

\MapSubheading{Symbolic comparison with Z3}
Z3 is a satisfiability modulo theories (SMT) solver: it searches for values that
make a collection of logical and arithmetic conditions true. Here the condition
asks whether two supported rule encodings return different answers. The facts,
output meanings and permitted input domain must first be made comparable.

\begin{MapDiagram}
\begin{tikzpicture}[node distance=8mm and 5mm]
\node[mapbox] (a) {Comparable rules\\Supported input domain};
\node[mapbox,right=of a] (b) {Encode a disagreement\\Ask Z3 for a case};
\node[mapbox,right=of b] (c) {Different answers found\\Replay in Python and Java};
\node[mapbox,below=of a] (d) {Unsupported, unknown\\or inconsistent domain};
\node[mapbox,below=of b] (e) {No counterexample\\Within the stated domain};
\draw[maparr](a)--(b);\draw[maparr](b)--(c);
\draw[maparr](b)--node[right,mapnote]{ruled out}(e);
\draw[maparr](a)--(d);
\draw[maparr](b.south west)--(d.north east);
\end{tikzpicture}
\end{MapDiagram}
\MapCaption{A solver result concerns an encoded comparison. An unsupported
operation, timeout or empty domain cannot count as agreement.}

The current adapter supports a restricted subset of RuleIR and does not check
an independent proof certificate from the solver. A found discrepancy is
valuable because the original programs can replay it. A no-counterexample
result retains the encoding and domain assumptions. Chapter~\ref{ch:proof}
explains those restrictions.

\MapSubheading{Proofs of selected mathematical properties}
Lean is a proof assistant that checks a formal derivation. The mathematical
review used it for selected logical and arithmetic propositions, including
three-valued logic, an exception counterexample and conservation under a stated
transfer. SymPy, a symbolic algebra system, checked selected algebraic identities.

\begin{MapDiagram}
\begin{tikzpicture}[node distance=5mm]
\node[mapbox] (a) {State a property\\Definitions and assumptions};
\node[mapbox,right=of a] (b) {Lean checks the proof\\Selected propositions};
\node[mapbox,right=of b] (c) {Retain the conclusion\\Only the property proved};
\draw[maparr](a)--(b);\draw[maparr](b)--(c);
\end{tikzpicture}
\end{MapDiagram}
\MapCaption{These are checks of selected mathematical claims, not an automatic
proof of each generated program or of the whole compiler.}

Proving that every supported rule preserves its behaviour through translation
remains a stronger goal. The present checks do not establish that whole-compiler
result or correct interpretation of English. The Stipula/JML/KeY route offers
another research approach to particular Java properties.

\clearpage
\MapHeading{7. Deliver a reviewable Java program}{map:delivery}
The result is a program that another person can inspect, run and question.
A Java archive (JAR) packages the compiled code. The associated review record
identifies the sources, proposed meanings, facts, test results and unresolved
questions. A separate host program exercises the callable Java control.

\begin{MapDiagram}
\begin{tikzpicture}[node distance=7mm and 5mm]
\node[mapbox] (a) {Compiled Java and tests\\Fixed proposed control};
\node[mapbox,right=of a] (b) {Sources and meaning\\Assumptions and questions};
\node[mapbox,right=of b] (c) {Reviewable draft\\JAR and host example};
\draw[maparr](a)--(b);\draw[maparr](b)--(c);
\end{tikzpicture}
\end{MapDiagram}
\MapCaption{The delivered evidence lets a reviewer relate an executable answer
to its proposed meaning. Draft delivery supplies no authority to use it.}

\MapHeading{8. Approve and supervise bank use}{map:operation}
The institution decides which interpretation and exact program version it will
authorize. Its own identities, fact owners and transaction systems must preserve
that decision. The local source monitor and bounded scheduler can revisit
affected controls; actual bank operation still needs institution validation.

\begin{MapDiagram}
\begin{tikzpicture}[node distance=9mm and 5mm]
\node[mapresearch] (a) {Bank review\\Meaning and facts\\Release approval};
\node[mapresearch,right=of a] (b) {Supervised bank use\\Host, identity and transaction checks};
\node[mapbox,right=of b] (c) {Source review\\Monitor and scheduler};
\node[mapbox,below=of b] (d) {Revisit affected rules\\Preserve earlier decisions};
\draw[maparr](a)--(b);\draw[maparr](b)--(c);
\draw[maparr](c.south)--++(0,-3mm)-|node[pos=.25,below,mapnote]{source changes}(d.north);
\draw[maparr](d.west)-|node[pos=.25,below,mapnote]{revised proposal}(a.south);
\end{tikzpicture}
\end{MapDiagram}
\MapCaption{Dashed boxes require the bank's operating evidence. Local monitoring
is implemented; an unchanged source still leaves unresolved questions open.}

\MapSubheading{Measuring whether the method helps}
Empirical evaluation runs alongside development. The implemented comparison
offers a single reader, isolated readers, tree search and source assurance under
declared resource limits. Independent reference outcomes and unfamiliar source
families are needed to estimate interpretation quality; a separate workflow
study measures people's review effort.

\begin{MapDiagram}
\begin{tikzpicture}[node distance=5mm]
\node[mapbox] (a) {Fixed study design\\Hidden reference answers};
\node[mapbox,right=of a] (b) {Compare methods\\Errors, useful answers\\and missing answers};
\node[mapbox,right=of b] (c) {Estimate performance\\By source family};
\draw[maparr](a)--(b);\draw[maparr](b)--(c);
\end{tikzpicture}
\end{MapDiagram}
\MapCaption{The comparison machinery exists. Its development checks have not
established a method ranking or a reduction in reviewer workload.}

Chapters~\ref{ch:implementation}--\ref{ch:operation} explain delivery,
comparative evaluation and bank adoption.
\endgroup
\clearpage
